
\vspace{-0.25cm}
% Dans la partie « Dynamique dans un référentiel non galiléen », l'étude du champ de pesanteur est conduite en supposant le référentiel géocentrique galiléen. De nombreuses applications permettent d'illustrer cette partie : le pendule de Foucault, la déviation vers l'est, les vents géostrophiques, les courants marins ; l'étude statique des marées constitue également une ouverture pertinente.
{\footnotesize
\begin{tabularx}{\textwidth}[t]{|X|X|}
  \hline 
  \textbf{Notions et contenus} & \textbf{Capacités exigibles} \\
  \hline
  
  Contraintes tangentielles dans un écoulement $v=v_x(y) u_x$ au sein d'un fluide newtonien ; viscosité. & Utiliser l'expression fournie  $d F=\eta \partial v_x / \partial y d S u_x .$ \\
  \hline 
  Équivalent volumique des forces de viscosité dans un écoulement incompressible. & Établir l'expression de l'équivalent volumique des forces de viscosité dans le cas d'un écoulement de cisaillement à une dimension et utiliser sa généralisation admise pour un écoulement incompressible quelconque. \\
  \hline 
  Traînée d'une sphère solide en mouvement rectiligne uniforme dans un fluide newtonien : nombre de Reynolds ; coefficient de traînée $\mathrm{C}_{\mathrm{x}}$; graphe de $\mathrm{C}_{\mathrm{x}}$ en fonction du nombre de Reynolds. & Évaluer un nombre de Reynolds pour choisir un modèle de traînée linéaire ou un modèle de traînée quadratique. \\
  \hline
  \multicolumn{2}{|l|}{4.3.3. Équations dynamiques locales } \\
  \hline 
  Équation de Navier-Stokes dans un fluide newtonien en écoulement incompressible. Terme convectif. Terme diffusif. Nombre de Reynolds dans le cas d'une unique échelle spatiale. & Équation de Navier-Stokes dans un fluide newtonien en écoulement incompressible. Terme convectif. Terme diffusif. Nombre de Reynolds dans le cas d'une unique échelle spatiale. \\ 
  \hline
\end{tabularx}
}